\chapter{Pipeline Integrado em Producao}
\label{chap:pipeline_integrado_producao}

\chapterepigraph{``Ideas are cheap. It's only what you do with them that counts.''}{Isaac Asimov, \textit{The Secrets of the Universe} (1989).}
{\small\itshape\color{AccentGray}
Asimov sabia que nenhuma ideia vale seu pre\c{c}o antes de enfrentar a execu\c{c}\~{a}o~\cite{asimov1989secrets}. Paulo Freire nomeou o problema inverso: o curr\'{i}culo fragmentado deposita conceitos isolados sem jamais reconstituir a totalidade concreta de que fazem parte --- e o resultado \'{e} um saber in\'{e}rte, incapaz de agir sobre o mundo~\cite{freire1968pedagogia}. Este cap\'{i}tulo \'{e} a resposta t\'{e}cnica a esse problema: os componentes constru\'{i}dos separadamente s\'{o} revelam seu sentido pleno quando operam simultaneamente, como sistema coeso. A integra\c{c}\~{a}o n\~{a}o \'{e} um ap\^{e}ndice --- \'{e} o momento em que as partes se tornam, finalmente, compreens\'{i}veis.\par}
\vspace{0.6em}
% ---------------------------------------------------------------
\section{Escopo do capitulo}
% ---------------------------------------------------------------

\begin{quote}\itshape
\textbf{Onde estamos na hist\'oria:} ponto de converg\^encia.
Os nove cap\'itulos anteriores constru\'iram pe\c{c}as separadas
--- o modelo, o agente, o RAG, o pipeline, os twins, o grafo.
Aqui elas se encaixam pela primeira vez em um sistema coeso.
Daqui em diante, a pergunta muda: n\~ao mais ``como construir?''
mas ``como garantir que funciona em produ\c{c}\~ao?''. O loop
externo come\c{c}a no pr\'oximo cap\'itulo.
\end{quote}

Este capitulo e o ponto de convergencia dos tres capitulos anteriores:
pipeline assincrono (Cap.~\ref{chap:pipeline_assincrono_aws}),
segmentacao e twins (Cap.~\ref{chap:segmentacao_personas_twins}) e
GraphRAG com Neptune (Cap.~\ref{chap:graphrag_neptune}).
Aqui os componentes sao observados juntos, como um sistema coeso em
producao, revelando decisoes de arquitetura que so fazem sentido quando
as tres camadas operam simultaneamente.

% ---------------------------------------------------------------
\section{Objetivos de aprendizagem}
% ---------------------------------------------------------------

Ao concluir este capitulo, o leitor sera capaz de:

\begin{itemize}[leftmargin=*]
  \item Distinguir as fases \textit{offline} (SageMaker Processing Job)
        e \textit{online} (Lambda por requisicao) do pipeline;
  \item Justificar por que a indexacao de twins usa indice unico com
        \texttt{pre\_filter} em vez de um indice por cliente;
  \item Calcular e interpretar o indice de confiabilidade a partir das
        fontes RAG ativas;
  \item Usar o \texttt{demo\_api.ipynb} como ferramenta de validacao
        ponta a ponta dos tres modos.
\end{itemize}

% ---------------------------------------------------------------
\section{Conceitos teoricos}
% ---------------------------------------------------------------

\subsection{Duas fases, um sistema}

O pipeline e dividido em duas fases com responsabilidades distintas:

\begin{itemize}[leftmargin=*]
  \item \textbf{Fase offline} (SageMaker Processing Job): execucao
        periodica ou por evento. Carrega dados brutos, treina
        K-Means\index{K-Means} com StandardScaler\index{StandardScaler},
        indexa os clusters no OpenSearch\index{OpenSearch}, indexa os
        digital twins\index{digital twin} e persiste o modelo em S3.
  \item \textbf{Fase online} (Lambda\index{Lambda (AWS)} por requisicao):
        no cold start, carrega o modelo do S3. Por requisicao, classifica
        o cliente, executa o agente LangGraph\index{LangGraph} com RAG
        hibrido e retorna a resposta com indice de confiabilidade.
\end{itemize}

\begin{lstlisting}[language=bash,
  caption={Visao esquematica das duas fases do pipeline.},
  label={lst:duas_fases}]
# FASE OFFLINE (SageMaker Job)
Dados brutos (S3)
  -> K-Means + StandardScaler
  -> Indexar OpenSearch: clientes-segmento-{0..3}
  -> Indexar OpenSearch: clientes-digital-twins
  -> Salvar modelo: s3://bucket/models/clustering.pkl

# FASE ONLINE (Lambda - cold start, uma vez por instancia)
carregar_pkl_s3("clustering.pkl")
BedrockEmbeddings(model="amazon.titan-embed-text-v1")
criar_agente_por_cluster(cluster_id)  # 1 agente ReAct por segmento

# FASE ONLINE (Lambda - por requisicao)
KMeans.predict(StandardScaler.transform(dados_cliente)) -> cluster_id
buscar_rag(cluster_id ou cliente_id)  # OpenSearch + Neptune
LangGraph ReAct -> resposta + indice_confiabilidade
\end{lstlisting}

\subsection{Estrategia de indexacao: indice unico versus indice por cliente}

O modo \texttt{twin}\index{digital twin} poderia usar um indice OpenSearch
por cliente --- solucao intuitiva, porem inviavel em producao:
$N$ clientes gerariam $N$ indices, elevando custo de gerenciamento,
tempo de cold start e complexidade de re-indexacao.

A estrategia adotada e um indice unico \texttt{clientes-digital-twins}
com o campo de metadado \texttt{cliente\_id}. Na consulta, aplica-se
\texttt{pre\_filter} para restringir o escore k-NN ao cliente especifico:

\begin{lstlisting}[language=Python,
  caption={Retriever com pre\_filter por cliente\_id no indice de twins.},
  label={lst:prefilter_twin}]
vs.as_retriever(
    search_kwargs={
        "k": 3,
        "pre_filter": {
            "term": {"metadata.cliente_id": cliente_id}
        },
    }
)
\end{lstlisting}

Essa abordagem mantem latencia equivalente a um unico indice e
permite re-indexacao limpa de clientes especificos sem afetar os demais.

\subsection{Indice de confiabilidade}

O indice de confiabilidade\index{indice de confiabilidade|textbf}
quantifica o quanto o LLM utilizou o contexto RAG recuperado.
E calculado em duas dimensoes:

\begin{equation}
  \text{score} = 0{,}85 \times \text{cobertura\_fontes}
               + 0{,}15 \times \text{sobreposicao\_lexica}
  \label{eq:indice_confiabilidade}
\end{equation}

A \textbf{cobertura de fontes} e a soma dos pesos das camadas RAG ativas,
conforme a Tabela~\ref{tab:pesos_rag}.

\begin{table}[ht]
\centering
\caption{Pesos das fontes RAG no indice de confiabilidade.}
\label{tab:pesos_rag}
\begin{tabular}{lcc}
\toprule
\textbf{Fonte} & \textbf{Peso} & \textbf{Condicao de ativacao} \\
\midrule
BM25 OpenSearch  & 0,40 & Ao menos 1 documento recuperado \\
Neptune live     & 0,30 & Query OpenCypher retorna resultados \\
Neptune sync     & 0,15 & Snapshot replicado disponivel \\
\midrule
\textbf{Maximo}  & \textbf{0,85} & Todas as fontes ativas \\
\bottomrule
\end{tabular}
\end{table}

A \textbf{sobreposicao lexica} mede a fracao de termos-chave do contexto
recuperado que aparece na resposta final, com contribuicao maxima de 0,15.

A escala de interpretacao do score final e:

\begin{itemize}[leftmargin=*]
  \item $\geq 0{,}70$ --- \textbf{ALTO}: duas ou mais fontes ativas;
        LLM utilizou o contexto de forma verificavel.
  \item $0{,}40$--$0{,}69$ --- \textbf{MEDIO}: uma fonte ativa ou
        sobreposicao lexica parcial.
  \item $< 0{,}40$ --- \textbf{BAIXO}: sem contexto recuperado ---
        resposta baseada em conhecimento parametrico do LLM.
\end{itemize}

\textbf{Nota sobre o modo persona:} a sobreposicao lexica tende a ser
menor porque o LLM reescreve o contexto em linguagem coloquial de
primeira pessoa, mas a cobertura de fontes reflete corretamente que
ambas as camadas foram consultadas.

% ---------------------------------------------------------------
\section{Implementacao orientada por passos}
% ---------------------------------------------------------------

\subsection{Passo 1 --- Executar a fase offline}

\begin{lstlisting}[language=bash,
  caption={Execucao do pipeline offline para gerar indices e modelo.},
  label={lst:pipeline_offline}]
# Em ambiente SageMaker ou local com credenciais AWS configuradas:
python aws_pipeline_clientes.py --mode pipeline

# Verificar artefato gerado no S3:
aws s3 ls s3://$S3_BUCKET/models/clustering.pkl

# Verificar indices criados no OpenSearch:
curl -s "$OS_ENDPOINT/_cat/indices?v&h=index,docs.count" \
  | grep clientes
\end{lstlisting}

Resultado esperado: quatro indices \texttt{clientes-segmento-\{0..3\}}
e um indice \texttt{clientes-digital-twins}, cada um com a contagem de
documentos correspondente ao numero de clientes indexados.

\subsection{Passo 2 --- Verificar artefatos no S3 e indices no OpenSearch}

Apos a fase offline, os seguintes artefatos devem existir:

\begin{itemize}[leftmargin=*]
  \item \texttt{s3://bucket/models/clustering.pkl} --- modelo KMeans +
        StandardScaler + perfis enriquecidos;
  \item Cinco indices OpenSearch: quatro de segmento e um de twins.
\end{itemize}

\subsection{Passo 3 --- Executar os tres modos via notebook}

O laboratorio oficial deste capitulo e o \texttt{demo\_api.ipynb}.
Execute as celulas na ordem para observar os tres modos em producao.
Os modos diferem no sistema prompt e na fonte de RAG,
conforme a Tabela~\ref{tab:modos_agente}.

\begin{table}[ht]
\centering
\caption{Comparativo dos tres modos de agente.}
\label{tab:modos_agente}
\begin{tabular}{llll}
\toprule
\textbf{Modo} & \textbf{Perspectiva} & \textbf{Indice RAG} & \textbf{Requisito} \\
\midrule
\texttt{segmento} & 3a pessoa (gerente)    & \texttt{clientes-segmento-\{id\}} & \texttt{dados\_cliente} \\
\texttt{persona}  & 1a pessoa (arquetipo)  & \texttt{clientes-segmento-\{id\}} & \texttt{cluster\_id} \\
\texttt{twin}     & 1a pessoa (cliente real) & \texttt{clientes-digital-twins} & \texttt{cliente\_id} \\
\bottomrule
\end{tabular}
\end{table}

\subsection{Passo 4 --- Analisar o indice de confiabilidade}

Apos executar os tres modos, a celula de tabela comparativa do notebook
consolida os scores. Espera-se:

\begin{itemize}[leftmargin=*]
  \item \texttt{segmento} e \texttt{twin} com cliente indexado:
        score ALTO ($\geq 0{,}70$);
  \item \texttt{persona}: score ALTO em cobertura de fontes,
        sobreposicao lexica inferior ao modo segmento;
  \item \texttt{twin} com cliente sem dados no indice:
        score BAIXO ($< 0{,}40$).
\end{itemize}

\subsection{Passo 5 --- Simular degradacao por ausencia de indexacao}

A celula ``Demonstracao de degradacao'' do notebook usa o cliente
hipotetico \texttt{DEMO-X999} (inexistente no indice) para demonstrar
que o indice de confiabilidade detecta corretamente a ausencia de
contexto RAG --- o score cai para o nivel BAIXO e a resposta passa a
depender exclusivamente do conhecimento parametrico do LLM.

\subsection{Passo 6 --- Resultado completo da aula}

Ao final, o aluno deve comprovar:

\begin{itemize}[leftmargin=*]
  \item Fase offline executada com artefatos persistidos no S3;
  \item Quatro indices de segmento e um de twins no OpenSearch;
  \item Tres modos retornando \texttt{COMPLETED} com score de
        confiabilidade coerente com as fontes ativas;
  \item Score do twin com cliente indexado superior ao twin sem
        cliente indexado.
\end{itemize}

% ---------------------------------------------------------------
\section{Plano de testes do capitulo}
% ---------------------------------------------------------------

\subsection{Teste funcional}

\begin{itemize}[leftmargin=*]
  \item \textbf{TC-01:} executar fase offline e verificar existencia
        de \texttt{clustering.pkl} no S3 e cinco indices no OpenSearch.
  \item \textbf{TC-02:} modo \texttt{segmento} retorna
        \texttt{COMPLETED} com score $\geq 0{,}70$.
  \item \textbf{TC-03:} modo \texttt{twin} com cliente indexado retorna
        score superior ao mesmo modo com cliente nao indexado.
\end{itemize}

\subsection{Teste de validacao em producao (2026-05-03)}

Apos a reconstrucao completa da infraestrutura por produtos, o pipeline
integrado foi validado ponta a ponta no ambiente \texttt{dev} com os
recursos AWS efetivamente provisionados. Observou-se:

\begin{itemize}[leftmargin=*]
  \item \textbf{TVP-01 (modo segmento):}
        \texttt{request\_id}
        \nolinkurl{0e73a7f9-84ee-4359-aff3-54bd019b3410}. \par
        Estado final: \texttt{COMPLETED}; tempo de processamento:
        $\approx 18$\,s; \texttt{cluster\_id = 3}, segmento
        \emph{Massa Est\'{a}vel}; resposta gerada pelo modelo
        \texttt{global.anthropic.claude-sonnet-4-5-20250929-v1:0}
        via Amazon Bedrock;
  \item \textbf{Indice de confiabilidade:} \texttt{score = 0.0},
        n\'{\i}vel \emph{baixo}, \texttt{fontes\_ativas = []}, refletindo
        que o indice OpenSearch ainda n\~ao possu\'{\i}a documentos
        de segmento indexados --- confirma o comportamento esperado de
        degrada\c{c}\~ao graciosamente descrito no Passo~5;
  \item Modelo de segmenta\c{c}\~ao carregado de
        \texttt{s3://langchain-agent-artifacts-dev/\allowbreak{}clientes-agente/\allowbreak{}modelo\_clustering\_slim.pkl}
        com \texttt{n\_clusters = 4}.
\end{itemize}

Esse resultado valida TC-01 (modelo presente no S3) e TC-02 (pipeline retorna
\texttt{COMPLETED} no modo \texttt{segmento}), confirmando que a integra\c{c}\~ao
das tr\^es camadas --- classifica\c{c}\~ao, RAG e Bedrock --- opera corretamente
em produ\c{c}\~ao. O score baixo de confiabilidade sinaliza que a etapa
subsequente natural \'{e} executar a fase offline (\texttt{aws\_pipeline\_clientes.py
--mode pipeline}) para popular os \'{\i}ndices OpenSearch e habilitar o RAG h\'{\i}brido.

\subsection{Teste de qualidade de resposta}

\begin{itemize}[leftmargin=*]
  \item \textbf{TQ-01:} modo \texttt{persona} mantem voz em primeira
        pessoa em 100\% das respostas dos quatro clusters.
  \item \textbf{TQ-02:} modo \texttt{twin} com cliente indexado cita
        dados numericos do perfil individual em pelo menos 80\% das
        respostas.
\end{itemize}

% ---------------------------------------------------------------
\section{Criterios de aceite}
% ---------------------------------------------------------------

\begin{itemize}[leftmargin=*]
  \item TC-01 a TC-03 aprovados.
  \item TQ-01 e TQ-02 dentro dos limites estabelecidos.
  \item Score de confiabilidade calculado e exibido corretamente
        para os tres modos no notebook.
\end{itemize}

% ---------------------------------------------------------------
\section{Espaco para expansao iterativa}
% ---------------------------------------------------------------

\begin{itemize}[leftmargin=*]
  \item \textbf{v10.1} --- Adicionar cache de contexto RAG no DynamoDB
        para reduzir latencia em consultas repetidas ao mesmo cliente.
  \item \textbf{v10.2} --- Implementar re-indexacao incremental de twins
        sem necessidade de re-executar o pipeline completo.
  \item \textbf{v10.3} --- Expandir o indice de confiabilidade com
        metrica de coerencia semantica (cosine similarity) entre
        contexto recuperado e resposta gerada.
\end{itemize}
